% Options for packages loaded elsewhere
\PassOptionsToPackage{unicode}{hyperref}
\PassOptionsToPackage{hyphens}{url}
\documentclass[
]{article}
\usepackage{xcolor}
\usepackage[margin=1in]{geometry}
\usepackage{amsmath,amssymb}
\setcounter{secnumdepth}{-\maxdimen} % remove section numbering
\usepackage{iftex}
\ifPDFTeX
  \usepackage[T1]{fontenc}
  \usepackage[utf8]{inputenc}
  \usepackage{textcomp} % provide euro and other symbols
\else % if luatex or xetex
  \usepackage{unicode-math} % this also loads fontspec
  \defaultfontfeatures{Scale=MatchLowercase}
  \defaultfontfeatures[\rmfamily]{Ligatures=TeX,Scale=1}
\fi
\usepackage{lmodern}
\ifPDFTeX\else
  % xetex/luatex font selection
\fi
% Use upquote if available, for straight quotes in verbatim environments
\IfFileExists{upquote.sty}{\usepackage{upquote}}{}
\IfFileExists{microtype.sty}{% use microtype if available
  \usepackage[]{microtype}
  \UseMicrotypeSet[protrusion]{basicmath} % disable protrusion for tt fonts
}{}
\makeatletter
\@ifundefined{KOMAClassName}{% if non-KOMA class
  \IfFileExists{parskip.sty}{%
    \usepackage{parskip}
  }{% else
    \setlength{\parindent}{0pt}
    \setlength{\parskip}{6pt plus 2pt minus 1pt}}
}{% if KOMA class
  \KOMAoptions{parskip=half}}
\makeatother
\usepackage{color}
\usepackage{fancyvrb}
\newcommand{\VerbBar}{|}
\newcommand{\VERB}{\Verb[commandchars=\\\{\}]}
\DefineVerbatimEnvironment{Highlighting}{Verbatim}{commandchars=\\\{\}}
% Add ',fontsize=\small' for more characters per line
\usepackage{framed}
\definecolor{shadecolor}{RGB}{248,248,248}
\newenvironment{Shaded}{\begin{snugshade}}{\end{snugshade}}
\newcommand{\AlertTok}[1]{\textcolor[rgb]{0.94,0.16,0.16}{#1}}
\newcommand{\AnnotationTok}[1]{\textcolor[rgb]{0.56,0.35,0.01}{\textbf{\textit{#1}}}}
\newcommand{\AttributeTok}[1]{\textcolor[rgb]{0.13,0.29,0.53}{#1}}
\newcommand{\BaseNTok}[1]{\textcolor[rgb]{0.00,0.00,0.81}{#1}}
\newcommand{\BuiltInTok}[1]{#1}
\newcommand{\CharTok}[1]{\textcolor[rgb]{0.31,0.60,0.02}{#1}}
\newcommand{\CommentTok}[1]{\textcolor[rgb]{0.56,0.35,0.01}{\textit{#1}}}
\newcommand{\CommentVarTok}[1]{\textcolor[rgb]{0.56,0.35,0.01}{\textbf{\textit{#1}}}}
\newcommand{\ConstantTok}[1]{\textcolor[rgb]{0.56,0.35,0.01}{#1}}
\newcommand{\ControlFlowTok}[1]{\textcolor[rgb]{0.13,0.29,0.53}{\textbf{#1}}}
\newcommand{\DataTypeTok}[1]{\textcolor[rgb]{0.13,0.29,0.53}{#1}}
\newcommand{\DecValTok}[1]{\textcolor[rgb]{0.00,0.00,0.81}{#1}}
\newcommand{\DocumentationTok}[1]{\textcolor[rgb]{0.56,0.35,0.01}{\textbf{\textit{#1}}}}
\newcommand{\ErrorTok}[1]{\textcolor[rgb]{0.64,0.00,0.00}{\textbf{#1}}}
\newcommand{\ExtensionTok}[1]{#1}
\newcommand{\FloatTok}[1]{\textcolor[rgb]{0.00,0.00,0.81}{#1}}
\newcommand{\FunctionTok}[1]{\textcolor[rgb]{0.13,0.29,0.53}{\textbf{#1}}}
\newcommand{\ImportTok}[1]{#1}
\newcommand{\InformationTok}[1]{\textcolor[rgb]{0.56,0.35,0.01}{\textbf{\textit{#1}}}}
\newcommand{\KeywordTok}[1]{\textcolor[rgb]{0.13,0.29,0.53}{\textbf{#1}}}
\newcommand{\NormalTok}[1]{#1}
\newcommand{\OperatorTok}[1]{\textcolor[rgb]{0.81,0.36,0.00}{\textbf{#1}}}
\newcommand{\OtherTok}[1]{\textcolor[rgb]{0.56,0.35,0.01}{#1}}
\newcommand{\PreprocessorTok}[1]{\textcolor[rgb]{0.56,0.35,0.01}{\textit{#1}}}
\newcommand{\RegionMarkerTok}[1]{#1}
\newcommand{\SpecialCharTok}[1]{\textcolor[rgb]{0.81,0.36,0.00}{\textbf{#1}}}
\newcommand{\SpecialStringTok}[1]{\textcolor[rgb]{0.31,0.60,0.02}{#1}}
\newcommand{\StringTok}[1]{\textcolor[rgb]{0.31,0.60,0.02}{#1}}
\newcommand{\VariableTok}[1]{\textcolor[rgb]{0.00,0.00,0.00}{#1}}
\newcommand{\VerbatimStringTok}[1]{\textcolor[rgb]{0.31,0.60,0.02}{#1}}
\newcommand{\WarningTok}[1]{\textcolor[rgb]{0.56,0.35,0.01}{\textbf{\textit{#1}}}}
\usepackage{graphicx}
\makeatletter
\newsavebox\pandoc@box
\newcommand*\pandocbounded[1]{% scales image to fit in text height/width
  \sbox\pandoc@box{#1}%
  \Gscale@div\@tempa{\textheight}{\dimexpr\ht\pandoc@box+\dp\pandoc@box\relax}%
  \Gscale@div\@tempb{\linewidth}{\wd\pandoc@box}%
  \ifdim\@tempb\p@<\@tempa\p@\let\@tempa\@tempb\fi% select the smaller of both
  \ifdim\@tempa\p@<\p@\scalebox{\@tempa}{\usebox\pandoc@box}%
  \else\usebox{\pandoc@box}%
  \fi%
}
% Set default figure placement to htbp
\def\fps@figure{htbp}
\makeatother
\setlength{\emergencystretch}{3em} % prevent overfull lines
\providecommand{\tightlist}{%
  \setlength{\itemsep}{0pt}\setlength{\parskip}{0pt}}
\usepackage{bookmark}
\IfFileExists{xurl.sty}{\usepackage{xurl}}{} % add URL line breaks if available
\urlstyle{same}
\hypersetup{
  pdftitle={Lab 2 - R Basics, Data Types, Data Manipulation},
  hidelinks,
  pdfcreator={LaTeX via pandoc}}

\title{Lab 2 - R Basics, Data Types, Data Manipulation}
\author{}
\date{\vspace{-2.5em}2026-01-16}

\begin{document}
\maketitle

\subsection{Learning Objectives}\label{learning-objectives}

After completing this lab, you will be able to:

\begin{enumerate}
\def\labelenumi{\arabic{enumi}.}
\tightlist
\item
  Write clear, well-commented code
\item
  Understand and use variables to store data
\item
  Call functions and provide arguments
\item
  Work with different data types (numeric, character, logical)
\item
  Create and manipulate vectors and data frames
\item
  Install and load R packages
\end{enumerate}

\subsection{How to Use This Script}\label{how-to-use-this-script}

\begin{enumerate}
\def\labelenumi{\arabic{enumi}.}
\tightlist
\item
  Read through the explanatory text
\item
  Run each code block one at a time (click the green play button or
  press Ctrl+Enter / Cmd+Enter)
\item
  Observe what happens in the output below each chunk
\item
  Complete the exercises marked ``EXERCISE''
\end{enumerate}

\begin{center}\rule{0.5\linewidth}{0.5pt}\end{center}

\section{PART 1: Comments - Documenting Your
Code}\label{part-1-comments---documenting-your-code}

Comments start with the \texttt{\#} symbol. R ignores everything after
\texttt{\#} on a line.

\textbf{Why use comments?}

\begin{itemize}
\tightlist
\item
  Explain what your code does (for yourself and others)
\item
  Leave notes for future reference
\item
  Temporarily disable code without deleting it
\end{itemize}

\begin{Shaded}
\begin{Highlighting}[]
\CommentTok{\# Bad comment:}
\NormalTok{n }\OtherTok{\textless{}{-}} \DecValTok{5}  \CommentTok{\# set n to 5}

\CommentTok{\# Good comment:}
\NormalTok{n }\OtherTok{\textless{}{-}} \DecValTok{5}  \CommentTok{\# set sample size }
\end{Highlighting}
\end{Shaded}

\begin{center}\rule{0.5\linewidth}{0.5pt}\end{center}

\section{PART 2: Variables - Storing
Information}\label{part-2-variables---storing-information}

A \textbf{variable} is a labeled container that holds a value. Use
\texttt{\textless{}-} (assignment arrow) to store a value in a variable.
Keyboard shortcut: Alt and - (Windows) or Option and - (Mac).

\begin{Shaded}
\begin{Highlighting}[]
\NormalTok{sample\_size }\OtherTok{\textless{}{-}} \DecValTok{30}
\NormalTok{temperature\_celsius }\OtherTok{\textless{}{-}} \FloatTok{22.5}
\NormalTok{species\_name }\OtherTok{\textless{}{-}} \StringTok{"Carcharhinus melanopterus"}
\end{Highlighting}
\end{Shaded}

To see what's stored in a variable, just type its name:

\begin{Shaded}
\begin{Highlighting}[]
\NormalTok{sample\_size}
\end{Highlighting}
\end{Shaded}

\begin{verbatim}
## [1] 30
\end{verbatim}

\begin{Shaded}
\begin{Highlighting}[]
\NormalTok{temperature\_celsius}
\end{Highlighting}
\end{Shaded}

\begin{verbatim}
## [1] 22.5
\end{verbatim}

\begin{Shaded}
\begin{Highlighting}[]
\NormalTok{species\_name}
\end{Highlighting}
\end{Shaded}

\begin{verbatim}
## [1] "Carcharhinus melanopterus"
\end{verbatim}

You can use variables in calculations and store the output in a new
variable name:

\begin{Shaded}
\begin{Highlighting}[]
\NormalTok{temperature\_fahrenheit }\OtherTok{\textless{}{-}}\NormalTok{ (temperature\_celsius }\SpecialCharTok{*} \DecValTok{9}\SpecialCharTok{/}\DecValTok{5}\NormalTok{) }\SpecialCharTok{+} \DecValTok{32}
\NormalTok{temperature\_fahrenheit}
\end{Highlighting}
\end{Shaded}

\begin{verbatim}
## [1] 72.5
\end{verbatim}

Variables can be updated - this is a common source of error so keep good
track:

\begin{Shaded}
\begin{Highlighting}[]
\NormalTok{sample\_size }\OtherTok{\textless{}{-}} \DecValTok{50}
\NormalTok{sample\_size}
\end{Highlighting}
\end{Shaded}

\begin{verbatim}
## [1] 50
\end{verbatim}

Finally, take a look at your Environment pane to the right to see what
variables you have defined.

\subsubsection{Naming Rules for
Variables}\label{naming-rules-for-variables}

\begin{enumerate}
\def\labelenumi{\arabic{enumi}.}
\tightlist
\item
  Can contain letters, numbers, underscores (\texttt{\_}), and periods
  (\texttt{.})
\item
  Must START with a letter (or a period)
\item
  Case-sensitive: ``Species'' and ``species'' are different!
\item
  Cannot use reserved words (like: if, else, for, TRUE, FALSE)
\end{enumerate}

\begin{Shaded}
\begin{Highlighting}[]
\CommentTok{\# Good variable names:}
\NormalTok{body\_mass\_kg }\OtherTok{\textless{}{-}} \FloatTok{1.2}
\NormalTok{num\_observations }\OtherTok{\textless{}{-}} \DecValTok{100}
\NormalTok{is\_endangered }\OtherTok{\textless{}{-}} \ConstantTok{TRUE}

\CommentTok{\# Avoid these:}
\CommentTok{\# x \textless{}{-} 5           \# Not descriptive}
\CommentTok{\# 2nd\_sample \textless{}{-} 10 \# Can\textquotesingle{}t start with a number}
\CommentTok{\# my{-}data \textless{}{-} 5     \# Hyphens not allowed (use underscores)}
\end{Highlighting}
\end{Shaded}

\subsubsection{EXERCISE}\label{exercise}

Create a variable called \texttt{bird\_wingspan\_cm} and assign it a
value of 105. Then calculate the wingspan in meters and store it in
\texttt{bird\_wingspan\_m}.

\begin{Shaded}
\begin{Highlighting}[]
\CommentTok{\# Your code here:}
\NormalTok{bird\_wingspan\_cm }\OtherTok{\textless{}{-}} \DecValTok{105}
\NormalTok{bird\_wingspan\_m }\OtherTok{\textless{}{-}}\NormalTok{ bird\_wingspan\_cm}\SpecialCharTok{/}\DecValTok{100}
\end{Highlighting}
\end{Shaded}

\begin{center}\rule{0.5\linewidth}{0.5pt}\end{center}

\section{PART 3: Data Types}\label{part-3-data-types}

R works with several fundamental data types. When you store data in a
variable, R recognizes it as one of several types.

\subsection{1. NUMERIC - Numbers (integers and
decimals)}\label{numeric---numbers-integers-and-decimals}

\begin{Shaded}
\begin{Highlighting}[]
\NormalTok{population\_size }\OtherTok{\textless{}{-}} \DecValTok{1500}      \CommentTok{\# integer}
\NormalTok{average\_weight }\OtherTok{\textless{}{-}} \FloatTok{2.45}       \CommentTok{\# decimal (called "double")}
\NormalTok{temperature }\OtherTok{\textless{}{-}} \SpecialCharTok{{-}}\FloatTok{5.2}          \CommentTok{\# negative numbers work too}
\NormalTok{scientific }\OtherTok{\textless{}{-}} \FloatTok{6.022e23}       \CommentTok{\# scientific notation }

\CommentTok{\# Check the type with the class function:}
\FunctionTok{class}\NormalTok{(population\_size)  }\CommentTok{\# "numeric"}
\end{Highlighting}
\end{Shaded}

\begin{verbatim}
## [1] "numeric"
\end{verbatim}

\subsection{OPERATIONS WITH NUMERIC DATA
TYPES}\label{operations-with-numeric-data-types}

\begin{Shaded}
\begin{Highlighting}[]
\DecValTok{2} \SpecialCharTok{+} \DecValTok{1}   \CommentTok{\# Addition}
\end{Highlighting}
\end{Shaded}

\begin{verbatim}
## [1] 3
\end{verbatim}

\begin{Shaded}
\begin{Highlighting}[]
\DecValTok{2} \SpecialCharTok{{-}} \DecValTok{1}   \CommentTok{\# Subtraction }
\end{Highlighting}
\end{Shaded}

\begin{verbatim}
## [1] 1
\end{verbatim}

\begin{Shaded}
\begin{Highlighting}[]
\DecValTok{2} \SpecialCharTok{*} \DecValTok{1}   \CommentTok{\# Multiplication}
\end{Highlighting}
\end{Shaded}

\begin{verbatim}
## [1] 2
\end{verbatim}

\begin{Shaded}
\begin{Highlighting}[]
\DecValTok{2} \SpecialCharTok{/} \DecValTok{1}   \CommentTok{\# Division}
\end{Highlighting}
\end{Shaded}

\begin{verbatim}
## [1] 2
\end{verbatim}

\begin{Shaded}
\begin{Highlighting}[]
\DecValTok{2} \SpecialCharTok{\^{}} \DecValTok{2}   \CommentTok{\# Exponentiation}
\end{Highlighting}
\end{Shaded}

\begin{verbatim}
## [1] 4
\end{verbatim}

\begin{Shaded}
\begin{Highlighting}[]
\DecValTok{7} \SpecialCharTok{\%\%} \DecValTok{2}  \CommentTok{\# Modulo {-} remainder after division}
\end{Highlighting}
\end{Shaded}

\begin{verbatim}
## [1] 1
\end{verbatim}

\begin{Shaded}
\begin{Highlighting}[]
\DecValTok{7} \SpecialCharTok{\%/\%} \DecValTok{2} \CommentTok{\# Integer division}
\end{Highlighting}
\end{Shaded}

\begin{verbatim}
## [1] 3
\end{verbatim}

\subsection{2. CHARACTER - Text (also called
``strings'')}\label{character---text-also-called-strings}

Always use quotes (single or double).

\begin{Shaded}
\begin{Highlighting}[]
\NormalTok{genus }\OtherTok{\textless{}{-}} \StringTok{"Carcharhinus"}
\NormalTok{species }\OtherTok{\textless{}{-}} \StringTok{\textquotesingle{}melanopterus\textquotesingle{}}
\NormalTok{full\_name }\OtherTok{\textless{}{-}} \StringTok{"Carcharhinus melanopterus"}
\NormalTok{full\_name2 }\OtherTok{\textless{}{-}} \FunctionTok{paste}\NormalTok{(genus, species)}
\NormalTok{full\_name2}
\end{Highlighting}
\end{Shaded}

\begin{verbatim}
## [1] "Carcharhinus melanopterus"
\end{verbatim}

\begin{Shaded}
\begin{Highlighting}[]
\CommentTok{\# Check the type:}
\FunctionTok{class}\NormalTok{(genus)  }\CommentTok{\# "character"}
\end{Highlighting}
\end{Shaded}

\begin{verbatim}
## [1] "character"
\end{verbatim}

\textbf{IMPORTANT:} Numbers in quotes are TEXT, not numbers!

\begin{Shaded}
\begin{Highlighting}[]
\NormalTok{sample\_id }\OtherTok{\textless{}{-}} \StringTok{"42"} \CommentTok{\#numbers in quotations are characters"}
\FunctionTok{class}\NormalTok{(sample\_id)  }\CommentTok{\# "character" {-} we can\textquotesingle{}t perform operations with the variable sample\_id}
\end{Highlighting}
\end{Shaded}

\begin{verbatim}
## [1] "character"
\end{verbatim}

\begin{Shaded}
\begin{Highlighting}[]
\CommentTok{\#sample\_id + 1  \# generates an error because we\textquotesingle{}re trying to add an integer and a character}
\end{Highlighting}
\end{Shaded}

\subsection{3. LOGICAL - TRUE or FALSE (Boolean
values)}\label{logical---true-or-false-boolean-values}

Used for yes/no questions and conditions.

\begin{Shaded}
\begin{Highlighting}[]
\NormalTok{is\_mammal }\OtherTok{\textless{}{-}} \ConstantTok{FALSE}
\NormalTok{has\_wings }\OtherTok{\textless{}{-}} \ConstantTok{TRUE}
\NormalTok{is\_extinct }\OtherTok{\textless{}{-}} \ConstantTok{FALSE}

\FunctionTok{class}\NormalTok{(is\_mammal)  }\CommentTok{\# "logical"}
\end{Highlighting}
\end{Shaded}

\begin{verbatim}
## [1] "logical"
\end{verbatim}

Logical values come from comparisons:

\begin{Shaded}
\begin{Highlighting}[]
\DecValTok{5} \SpecialCharTok{\textgreater{}} \DecValTok{3}         \CommentTok{\# TRUE (5 is greater than 3) \#asking R if 5 is greater than 3, will respond "5"}
\end{Highlighting}
\end{Shaded}

\begin{verbatim}
## [1] TRUE
\end{verbatim}

\begin{Shaded}
\begin{Highlighting}[]
\DecValTok{5} \SpecialCharTok{\textless{}} \DecValTok{3}         \CommentTok{\# FALSE}
\end{Highlighting}
\end{Shaded}

\begin{verbatim}
## [1] FALSE
\end{verbatim}

\begin{Shaded}
\begin{Highlighting}[]
\DecValTok{5} \SpecialCharTok{==} \DecValTok{5} \CommentTok{\# TRUE (note: double equals for comparison)}
\end{Highlighting}
\end{Shaded}

\begin{verbatim}
## [1] TRUE
\end{verbatim}

\begin{Shaded}
\begin{Highlighting}[]
\DecValTok{5} \SpecialCharTok{!=} \DecValTok{3}        \CommentTok{\# TRUE (5 is NOT equal to 3)}
\end{Highlighting}
\end{Shaded}

\begin{verbatim}
## [1] TRUE
\end{verbatim}

\begin{Shaded}
\begin{Highlighting}[]
\DecValTok{10} \SpecialCharTok{\textgreater{}=} \DecValTok{10}      \CommentTok{\# TRUE (greater than OR equal)}
\end{Highlighting}
\end{Shaded}

\begin{verbatim}
## [1] TRUE
\end{verbatim}

\begin{Shaded}
\begin{Highlighting}[]
\DecValTok{10} \SpecialCharTok{\textless{}=} \DecValTok{5}       \CommentTok{\# FALSE}
\end{Highlighting}
\end{Shaded}

\begin{verbatim}
## [1] FALSE
\end{verbatim}

Logical operators:

\begin{Shaded}
\begin{Highlighting}[]
\ConstantTok{TRUE} \SpecialCharTok{\&} \ConstantTok{TRUE}   \CommentTok{\# AND: both must be true {-}\textgreater{} TRUE}
\end{Highlighting}
\end{Shaded}

\begin{verbatim}
## [1] TRUE
\end{verbatim}

\begin{Shaded}
\begin{Highlighting}[]
\ConstantTok{TRUE} \SpecialCharTok{\&} \ConstantTok{FALSE}  \CommentTok{\# {-}\textgreater{} FALSE}
\end{Highlighting}
\end{Shaded}

\begin{verbatim}
## [1] FALSE
\end{verbatim}

\begin{Shaded}
\begin{Highlighting}[]
\ConstantTok{TRUE} \SpecialCharTok{|} \ConstantTok{FALSE}  \CommentTok{\# OR: at least one must be true {-}\textgreater{} TRUE}
\end{Highlighting}
\end{Shaded}

\begin{verbatim}
## [1] TRUE
\end{verbatim}

\begin{Shaded}
\begin{Highlighting}[]
\SpecialCharTok{!}\ConstantTok{TRUE}         \CommentTok{\# NOT: flips the value {-}\textgreater{} FALSE}
\end{Highlighting}
\end{Shaded}

\begin{verbatim}
## [1] FALSE
\end{verbatim}

\begin{Shaded}
\begin{Highlighting}[]
\NormalTok{is\_mammal }\OtherTok{\textless{}{-}} \ConstantTok{TRUE}
\NormalTok{can\_fly }\OtherTok{\textless{}{-}} \ConstantTok{FALSE}
\NormalTok{is\_mammal }\SpecialCharTok{\&}\NormalTok{ can\_fly }\CommentTok{\#cannot be a mammal,, and fly, these values were prev assigned }
\end{Highlighting}
\end{Shaded}

\begin{verbatim}
## [1] FALSE
\end{verbatim}

\begin{Shaded}
\begin{Highlighting}[]
\NormalTok{is\_mammal }\SpecialCharTok{|}\NormalTok{ can\_fly}
\end{Highlighting}
\end{Shaded}

\begin{verbatim}
## [1] TRUE
\end{verbatim}

\subsection{4. SPECIAL VALUES}\label{special-values}

\begin{Shaded}
\begin{Highlighting}[]
\CommentTok{\# NA {-} "Not Available" (missing data {-} VERY common in biology!)}
\NormalTok{weight }\OtherTok{\textless{}{-}} \ConstantTok{NA}
\FunctionTok{is.na}\NormalTok{(weight)  }\CommentTok{\# TRUE {-} check if something is NA}
\end{Highlighting}
\end{Shaded}

\begin{verbatim}
## [1] TRUE
\end{verbatim}

\begin{Shaded}
\begin{Highlighting}[]
\CommentTok{\# NULL {-} empty (no value at all)}
\NormalTok{empty }\OtherTok{\textless{}{-}} \ConstantTok{NULL}
\FunctionTok{is.null}\NormalTok{(empty)  }\CommentTok{\# TRUE}
\end{Highlighting}
\end{Shaded}

\begin{verbatim}
## [1] TRUE
\end{verbatim}

\begin{Shaded}
\begin{Highlighting}[]
\CommentTok{\# NaN {-} "Not a Number" (undefined mathematical result)}
\DecValTok{0}\SpecialCharTok{/}\DecValTok{0}  \CommentTok{\# NaN}
\end{Highlighting}
\end{Shaded}

\begin{verbatim}
## [1] NaN
\end{verbatim}

\begin{Shaded}
\begin{Highlighting}[]
\CommentTok{\# Inf {-} Infinity}
\DecValTok{1}\SpecialCharTok{/}\DecValTok{0}   \CommentTok{\# Inf}
\end{Highlighting}
\end{Shaded}

\begin{verbatim}
## [1] Inf
\end{verbatim}

\begin{Shaded}
\begin{Highlighting}[]
\SpecialCharTok{{-}}\DecValTok{1}\SpecialCharTok{/}\DecValTok{0}  \CommentTok{\# {-}Inf}
\end{Highlighting}
\end{Shaded}

\begin{verbatim}
## [1] -Inf
\end{verbatim}

\subsection{Type Conversion}\label{type-conversion}

Sometimes you need to convert between types:

\begin{Shaded}
\begin{Highlighting}[]
\CommentTok{\# Text to number:}
\FunctionTok{as.numeric}\NormalTok{(}\StringTok{"42"}\NormalTok{)     }\CommentTok{\# 42}
\end{Highlighting}
\end{Shaded}

\begin{verbatim}
## [1] 42
\end{verbatim}

\begin{Shaded}
\begin{Highlighting}[]
\CommentTok{\# Number to text:}
\FunctionTok{as.character}\NormalTok{(}\DecValTok{42}\NormalTok{)     }\CommentTok{\# "42"}
\end{Highlighting}
\end{Shaded}

\begin{verbatim}
## [1] "42"
\end{verbatim}

\begin{Shaded}
\begin{Highlighting}[]
\CommentTok{\# To logical:}
\FunctionTok{as.logical}\NormalTok{(}\DecValTok{1}\NormalTok{)        }\CommentTok{\# TRUE (any non{-}zero is TRUE)}
\end{Highlighting}
\end{Shaded}

\begin{verbatim}
## [1] TRUE
\end{verbatim}

\begin{Shaded}
\begin{Highlighting}[]
\FunctionTok{as.logical}\NormalTok{(}\DecValTok{0}\NormalTok{)        }\CommentTok{\# FALSE}
\end{Highlighting}
\end{Shaded}

\begin{verbatim}
## [1] FALSE
\end{verbatim}

\begin{Shaded}
\begin{Highlighting}[]
\CommentTok{\# Invalid conversions give NA:}
\FunctionTok{as.numeric}\NormalTok{(}\StringTok{"hello"}\NormalTok{)  }\CommentTok{\# NA (with a warning)}
\end{Highlighting}
\end{Shaded}

\begin{verbatim}
## [1] NA
\end{verbatim}

\subsubsection{EXERCISE}\label{exercise-1}

\begin{enumerate}
\def\labelenumi{\arabic{enumi}.}
\tightlist
\item
  Create a variable for your favorite animal's scientific name
  (character)
\item
  Create a variable for its average lifespan in years (numeric)
\item
  Create a logical variable indicating whether it's extinct or extant
\item
  Check the class of each variable
\item
  Convert your numeric variable to a character
\end{enumerate}

\begin{Shaded}
\begin{Highlighting}[]
\CommentTok{\# Your code here:}
\NormalTok{cat\_scientific\_name }\OtherTok{\textless{}{-}} \StringTok{"Felis catus"}
\NormalTok{cat\_ave\_lifespan\_years }\OtherTok{\textless{}{-}} \DecValTok{20} 
\NormalTok{is\_cat\_extinct }\OtherTok{\textless{}{-}} \ConstantTok{FALSE}

\FunctionTok{class}\NormalTok{(cat\_scientific\_name)}
\end{Highlighting}
\end{Shaded}

\begin{verbatim}
## [1] "character"
\end{verbatim}

\begin{Shaded}
\begin{Highlighting}[]
\FunctionTok{class}\NormalTok{(cat\_ave\_lifespan\_years)}
\end{Highlighting}
\end{Shaded}

\begin{verbatim}
## [1] "numeric"
\end{verbatim}

\begin{Shaded}
\begin{Highlighting}[]
\FunctionTok{class}\NormalTok{(is\_cat\_extinct)}
\end{Highlighting}
\end{Shaded}

\begin{verbatim}
## [1] "logical"
\end{verbatim}

\begin{Shaded}
\begin{Highlighting}[]
\NormalTok{ave\_lifespan }\OtherTok{\textless{}{-}} \FunctionTok{as.character}\NormalTok{(cat\_ave\_lifespan\_years)}
\FunctionTok{is.character}\NormalTok{(ave\_lifespan)}
\end{Highlighting}
\end{Shaded}

\begin{verbatim}
## [1] TRUE
\end{verbatim}

\begin{center}\rule{0.5\linewidth}{0.5pt}\end{center}

\section{PART 4: Functions}\label{part-4-functions}

A \textbf{function} is a command that performs a specific task - it
often streamlines repetitive or complex tasks.

Anatomy of a function call:
\texttt{function\_name(argument1,\ argument2,\ ...)}

\textbf{Arguments} (or parameters) are the inputs you give to a
function.

\begin{Shaded}
\begin{Highlighting}[]
\CommentTok{\# Example: the sqrt() function calculates square roots}
\FunctionTok{sqrt}\NormalTok{(}\DecValTok{16}\NormalTok{)      }\CommentTok{\# Returns 4}
\end{Highlighting}
\end{Shaded}

\begin{verbatim}
## [1] 4
\end{verbatim}

\begin{Shaded}
\begin{Highlighting}[]
\FunctionTok{sqrt}\NormalTok{(}\DecValTok{2}\NormalTok{)       }\CommentTok{\# Returns 1.414...}
\end{Highlighting}
\end{Shaded}

\begin{verbatim}
## [1] 1.414214
\end{verbatim}

\begin{Shaded}
\begin{Highlighting}[]
\CommentTok{\# Example: the round() function rounds numbers}
\FunctionTok{round}\NormalTok{(}\FloatTok{3.14159}\NormalTok{)           }\CommentTok{\# Returns 3 (default: 0 decimal places)}
\end{Highlighting}
\end{Shaded}

\begin{verbatim}
## [1] 3
\end{verbatim}

\begin{Shaded}
\begin{Highlighting}[]
\FunctionTok{round}\NormalTok{(}\FloatTok{3.14159}\NormalTok{, }\AttributeTok{digits=}\DecValTok{2}\NormalTok{) }\CommentTok{\# Returns 3.14 (2 decimal places)}
\end{Highlighting}
\end{Shaded}

\begin{verbatim}
## [1] 3.14
\end{verbatim}

\begin{Shaded}
\begin{Highlighting}[]
\FunctionTok{round}\NormalTok{(}\FloatTok{3.14159}\NormalTok{, }\DecValTok{2}\NormalTok{)        }\CommentTok{\# Same thing (positional argument)}
\end{Highlighting}
\end{Shaded}

\begin{verbatim}
## [1] 3.14
\end{verbatim}

\begin{Shaded}
\begin{Highlighting}[]
\CommentTok{\# Example: the paste() function combines text}
\FunctionTok{paste}\NormalTok{(}\StringTok{"Hello"}\NormalTok{, }\StringTok{"World"}\NormalTok{)                    }\CommentTok{\# "Hello World"}
\end{Highlighting}
\end{Shaded}

\begin{verbatim}
## [1] "Hello World"
\end{verbatim}

\begin{Shaded}
\begin{Highlighting}[]
\FunctionTok{paste}\NormalTok{(}\StringTok{"Sample"}\NormalTok{, }\DecValTok{1}\NormalTok{, }\AttributeTok{sep=}\StringTok{"\_"}\NormalTok{)                }\CommentTok{\# "Sample\_1"}
\end{Highlighting}
\end{Shaded}

\begin{verbatim}
## [1] "Sample_1"
\end{verbatim}

\begin{Shaded}
\begin{Highlighting}[]
\FunctionTok{paste}\NormalTok{(}\StringTok{"A"}\NormalTok{, }\StringTok{"B"}\NormalTok{, }\StringTok{"C"}\NormalTok{, }\AttributeTok{sep=}\StringTok{"{-}"}\NormalTok{)              }\CommentTok{\# "A{-}B{-}C"}
\end{Highlighting}
\end{Shaded}

\begin{verbatim}
## [1] "A-B-C"
\end{verbatim}

\subsection{Getting Help}\label{getting-help}

\begin{Shaded}
\begin{Highlighting}[]
\NormalTok{?sqrt              }\CommentTok{\# Help page for sqrt function}
\FunctionTok{help}\NormalTok{(round)        }\CommentTok{\# Same as ?round}
\end{Highlighting}
\end{Shaded}

Google is your friend: ``R how to calculate standard deviation''

\subsection{Useful Built-in Functions}\label{useful-built-in-functions}

\begin{Shaded}
\begin{Highlighting}[]
\CommentTok{\# Math functions:}
\FunctionTok{abs}\NormalTok{(}\SpecialCharTok{{-}}\DecValTok{5}\NormalTok{)            }\CommentTok{\# Absolute value: 5}
\end{Highlighting}
\end{Shaded}

\begin{verbatim}
## [1] 5
\end{verbatim}

\begin{Shaded}
\begin{Highlighting}[]
\FunctionTok{sqrt}\NormalTok{(}\DecValTok{16}\NormalTok{)           }\CommentTok{\# Square root: 4}
\end{Highlighting}
\end{Shaded}

\begin{verbatim}
## [1] 4
\end{verbatim}

\begin{Shaded}
\begin{Highlighting}[]
\FunctionTok{log}\NormalTok{(}\DecValTok{10}\NormalTok{)            }\CommentTok{\# Natural log: 2.303}
\end{Highlighting}
\end{Shaded}

\begin{verbatim}
## [1] 2.302585
\end{verbatim}

\begin{Shaded}
\begin{Highlighting}[]
\FunctionTok{log10}\NormalTok{(}\DecValTok{100}\NormalTok{)         }\CommentTok{\# Log base 10: 2}
\end{Highlighting}
\end{Shaded}

\begin{verbatim}
## [1] 2
\end{verbatim}

\begin{Shaded}
\begin{Highlighting}[]
\FunctionTok{exp}\NormalTok{(}\DecValTok{1}\NormalTok{)             }\CommentTok{\# e\^{}1: 2.718}
\end{Highlighting}
\end{Shaded}

\begin{verbatim}
## [1] 2.718282
\end{verbatim}

\begin{Shaded}
\begin{Highlighting}[]
\FunctionTok{ceiling}\NormalTok{(}\FloatTok{3.2}\NormalTok{)       }\CommentTok{\# Round up: 4}
\end{Highlighting}
\end{Shaded}

\begin{verbatim}
## [1] 4
\end{verbatim}

\begin{Shaded}
\begin{Highlighting}[]
\FunctionTok{floor}\NormalTok{(}\FloatTok{3.9}\NormalTok{)         }\CommentTok{\# Round down: 3}
\end{Highlighting}
\end{Shaded}

\begin{verbatim}
## [1] 3
\end{verbatim}

\begin{Shaded}
\begin{Highlighting}[]
\CommentTok{\# Character functions:}
\FunctionTok{nchar}\NormalTok{(}\StringTok{"Biology"}\NormalTok{)           }\CommentTok{\# Number of characters: 7}
\end{Highlighting}
\end{Shaded}

\begin{verbatim}
## [1] 7
\end{verbatim}

\begin{Shaded}
\begin{Highlighting}[]
\FunctionTok{toupper}\NormalTok{(}\StringTok{"hello"}\NormalTok{)           }\CommentTok{\# "HELLO"}
\end{Highlighting}
\end{Shaded}

\begin{verbatim}
## [1] "HELLO"
\end{verbatim}

\begin{Shaded}
\begin{Highlighting}[]
\FunctionTok{tolower}\NormalTok{(}\StringTok{"WORLD"}\NormalTok{)           }\CommentTok{\# "world"}
\end{Highlighting}
\end{Shaded}

\begin{verbatim}
## [1] "world"
\end{verbatim}

\begin{Shaded}
\begin{Highlighting}[]
\FunctionTok{substr}\NormalTok{(}\StringTok{"Accipiter"}\NormalTok{, }\DecValTok{1}\NormalTok{, }\DecValTok{3}\NormalTok{)  }\CommentTok{\# Extract characters: "Acc"}
\end{Highlighting}
\end{Shaded}

\begin{verbatim}
## [1] "Acc"
\end{verbatim}

\subsubsection{EXERCISE}\label{exercise-2}

\begin{enumerate}
\def\labelenumi{\arabic{enumi}.}
\tightlist
\item
  Use \texttt{sqrt()} to calculate the square root of 144
\item
  Use \texttt{round()} to round 17.389 to 1 decimal place
\item
  Use \texttt{paste()} to combine your first and last name
\item
  Use \texttt{nchar()} to count the letters in ``Blacktip reef shark''
\end{enumerate}

\begin{Shaded}
\begin{Highlighting}[]
\CommentTok{\# Your code here:}
\FunctionTok{sqrt}\NormalTok{(}\DecValTok{144}\NormalTok{)}
\end{Highlighting}
\end{Shaded}

\begin{verbatim}
## [1] 12
\end{verbatim}

\begin{Shaded}
\begin{Highlighting}[]
\FunctionTok{round}\NormalTok{(}\FloatTok{17.389}\NormalTok{, }\DecValTok{1}\NormalTok{)}
\end{Highlighting}
\end{Shaded}

\begin{verbatim}
## [1] 17.4
\end{verbatim}

\begin{Shaded}
\begin{Highlighting}[]
\FunctionTok{paste}\NormalTok{(}\StringTok{"Tabitha"}\NormalTok{, }\StringTok{"Hafenbrak"}\NormalTok{, }\AttributeTok{sep=}\StringTok{"\_"}\NormalTok{)}
\end{Highlighting}
\end{Shaded}

\begin{verbatim}
## [1] "Tabitha_Hafenbrak"
\end{verbatim}

\begin{Shaded}
\begin{Highlighting}[]
\FunctionTok{nchar}\NormalTok{(}\StringTok{"Blacktip reef shark"}\NormalTok{)}
\end{Highlighting}
\end{Shaded}

\begin{verbatim}
## [1] 19
\end{verbatim}

\begin{center}\rule{0.5\linewidth}{0.5pt}\end{center}

\section{PART 5: Vectors}\label{part-5-vectors}

A \textbf{vector} is a sequence of values of the SAME type. It's the
fundamental data structure in R.

\subsection{Creating Vectors}\label{creating-vectors}

Use \texttt{c()} - which stands for ``combine'' or ``concatenate''.

\begin{Shaded}
\begin{Highlighting}[]
\CommentTok{\# Numeric vector (e.g., body masses of 5 birds in grams):}
\NormalTok{body\_mass }\OtherTok{\textless{}{-}} \FunctionTok{c}\NormalTok{(}\DecValTok{1200}\NormalTok{, }\DecValTok{1350}\NormalTok{, }\DecValTok{1150}\NormalTok{, }\DecValTok{1425}\NormalTok{, }\DecValTok{1275}\NormalTok{)}
\NormalTok{body\_mass}
\end{Highlighting}
\end{Shaded}

\begin{verbatim}
## [1] 1200 1350 1150 1425 1275
\end{verbatim}

\begin{Shaded}
\begin{Highlighting}[]
\CommentTok{\# Character vector (e.g., species names):}
\NormalTok{species }\OtherTok{\textless{}{-}} \FunctionTok{c}\NormalTok{(}\StringTok{"Goshawk"}\NormalTok{, }\StringTok{"Sparrowhawk"}\NormalTok{, }\StringTok{"Buzzard"}\NormalTok{, }\StringTok{"Kestrel"}\NormalTok{)}
\NormalTok{species}
\end{Highlighting}
\end{Shaded}

\begin{verbatim}
## [1] "Goshawk"     "Sparrowhawk" "Buzzard"     "Kestrel"
\end{verbatim}

\begin{Shaded}
\begin{Highlighting}[]
\CommentTok{\# Logical vector:}
\NormalTok{is\_raptor }\OtherTok{\textless{}{-}} \FunctionTok{c}\NormalTok{(}\ConstantTok{TRUE}\NormalTok{, }\ConstantTok{TRUE}\NormalTok{, }\ConstantTok{TRUE}\NormalTok{, }\ConstantTok{TRUE}\NormalTok{)}
\NormalTok{is\_raptor}
\end{Highlighting}
\end{Shaded}

\begin{verbatim}
## [1] TRUE TRUE TRUE TRUE
\end{verbatim}

\subsection{Shortcuts for Creating
Sequences}\label{shortcuts-for-creating-sequences}

\begin{Shaded}
\begin{Highlighting}[]
\CommentTok{\# The colon operator for integer sequences:}
\DecValTok{1}\SpecialCharTok{:}\DecValTok{10}              \CommentTok{\# 1, 2, 3, 4, 5, 6, 7, 8, 9, 10}
\end{Highlighting}
\end{Shaded}

\begin{verbatim}
##  [1]  1  2  3  4  5  6  7  8  9 10
\end{verbatim}

\begin{Shaded}
\begin{Highlighting}[]
\CommentTok{\# seq() for more control:}
\FunctionTok{seq}\NormalTok{(}\AttributeTok{from=}\DecValTok{0}\NormalTok{, }\AttributeTok{to=}\DecValTok{100}\NormalTok{, }\AttributeTok{by=}\DecValTok{10}\NormalTok{)      }\CommentTok{\# 0, 10, 20, ..., 100}
\end{Highlighting}
\end{Shaded}

\begin{verbatim}
##  [1]   0  10  20  30  40  50  60  70  80  90 100
\end{verbatim}

\begin{Shaded}
\begin{Highlighting}[]
\FunctionTok{seq}\NormalTok{(}\AttributeTok{from=}\DecValTok{0}\NormalTok{, }\AttributeTok{to=}\DecValTok{1}\NormalTok{, }\AttributeTok{length.out=}\DecValTok{5}\NormalTok{) }\CommentTok{\# 5 evenly spaced values}
\end{Highlighting}
\end{Shaded}

\begin{verbatim}
## [1] 0.00 0.25 0.50 0.75 1.00
\end{verbatim}

\begin{Shaded}
\begin{Highlighting}[]
\CommentTok{\# rep() to repeat values:}
\FunctionTok{rep}\NormalTok{(}\DecValTok{0}\NormalTok{, }\AttributeTok{times=}\DecValTok{5}\NormalTok{)                 }\CommentTok{\# 0, 0, 0, 0, 0}
\end{Highlighting}
\end{Shaded}

\begin{verbatim}
## [1] 0 0 0 0 0
\end{verbatim}

\begin{Shaded}
\begin{Highlighting}[]
\FunctionTok{rep}\NormalTok{(}\FunctionTok{c}\NormalTok{(}\StringTok{"A"}\NormalTok{, }\StringTok{"B"}\NormalTok{), }\AttributeTok{times=}\DecValTok{3}\NormalTok{)       }\CommentTok{\# "A", "B", "A", "B", "A", "B"}
\end{Highlighting}
\end{Shaded}

\begin{verbatim}
## [1] "A" "B" "A" "B" "A" "B"
\end{verbatim}

\begin{Shaded}
\begin{Highlighting}[]
\FunctionTok{rep}\NormalTok{(}\FunctionTok{c}\NormalTok{(}\StringTok{"A"}\NormalTok{, }\StringTok{"B"}\NormalTok{), }\AttributeTok{each=}\DecValTok{3}\NormalTok{)        }\CommentTok{\# "A", "A", "A", "B", "B", "B"}
\end{Highlighting}
\end{Shaded}

\begin{verbatim}
## [1] "A" "A" "A" "B" "B" "B"
\end{verbatim}

\subsection{Vector Properties}\label{vector-properties}

\begin{Shaded}
\begin{Highlighting}[]
\FunctionTok{length}\NormalTok{(body\_mass)     }\CommentTok{\# How many elements? 5}
\end{Highlighting}
\end{Shaded}

\begin{verbatim}
## [1] 5
\end{verbatim}

\begin{Shaded}
\begin{Highlighting}[]
\FunctionTok{class}\NormalTok{(body\_mass)      }\CommentTok{\# What type? "numeric"}
\end{Highlighting}
\end{Shaded}

\begin{verbatim}
## [1] "numeric"
\end{verbatim}

\begin{Shaded}
\begin{Highlighting}[]
\FunctionTok{summary}\NormalTok{(body\_mass)    }\CommentTok{\# Quick statistical summary}
\end{Highlighting}
\end{Shaded}

\begin{verbatim}
##    Min. 1st Qu.  Median    Mean 3rd Qu.    Max. 
##    1150    1200    1275    1280    1350    1425
\end{verbatim}

\subsection{Accessing Vector Elements
(Indexing)}\label{accessing-vector-elements-indexing}

SUPER IMPORTANT!!

R uses 1-based indexing (first element is position 1). For Python users
- Python uses 0-based, so the first element is position 0.

\begin{Shaded}
\begin{Highlighting}[]
\NormalTok{body\_mass[}\DecValTok{1}\NormalTok{]          }\CommentTok{\# First element: 1200}
\end{Highlighting}
\end{Shaded}

\begin{verbatim}
## [1] 1200
\end{verbatim}

\begin{Shaded}
\begin{Highlighting}[]
\NormalTok{body\_mass[}\DecValTok{3}\NormalTok{]          }\CommentTok{\# Third element: 1150}
\end{Highlighting}
\end{Shaded}

\begin{verbatim}
## [1] 1150
\end{verbatim}

\begin{Shaded}
\begin{Highlighting}[]
\NormalTok{body\_mass[}\FunctionTok{c}\NormalTok{(}\DecValTok{1}\NormalTok{,}\DecValTok{3}\NormalTok{,}\DecValTok{5}\NormalTok{)]   }\CommentTok{\# Elements 1, 3, and 5}
\end{Highlighting}
\end{Shaded}

\begin{verbatim}
## [1] 1200 1150 1275
\end{verbatim}

\begin{Shaded}
\begin{Highlighting}[]
\CommentTok{\# Negative indices EXCLUDE elements:}
\NormalTok{body\_mass[}\SpecialCharTok{{-}}\DecValTok{1}\NormalTok{]         }\CommentTok{\# Everything except first element}
\end{Highlighting}
\end{Shaded}

\begin{verbatim}
## [1] 1350 1150 1425 1275
\end{verbatim}

\begin{Shaded}
\begin{Highlighting}[]
\NormalTok{body\_mass[}\SpecialCharTok{{-}}\FunctionTok{c}\NormalTok{(}\DecValTok{1}\NormalTok{,}\DecValTok{2}\NormalTok{)]    }\CommentTok{\# Everything except elements 1 and 2}
\end{Highlighting}
\end{Shaded}

\begin{verbatim}
## [1] 1150 1425 1275
\end{verbatim}

\begin{Shaded}
\begin{Highlighting}[]
\NormalTok{body\_mass[}\FunctionTok{length}\NormalTok{(body\_mass)] }
\end{Highlighting}
\end{Shaded}

\begin{verbatim}
## [1] 1275
\end{verbatim}

\begin{Shaded}
\begin{Highlighting}[]
\CommentTok{\# Get a range:}
\NormalTok{body\_mass[}\DecValTok{2}\SpecialCharTok{:}\DecValTok{4}\NormalTok{]        }\CommentTok{\# Elements 2 through 4}
\end{Highlighting}
\end{Shaded}

\begin{verbatim}
## [1] 1350 1150 1425
\end{verbatim}

\subsection{Named Vectors}\label{named-vectors}

Give names to elements for easier access:

\begin{Shaded}
\begin{Highlighting}[]
\NormalTok{wing\_length }\OtherTok{\textless{}{-}} \FunctionTok{c}\NormalTok{(}\AttributeTok{goshawk=}\DecValTok{58}\NormalTok{, }\AttributeTok{sparrowhawk=}\DecValTok{32}\NormalTok{, }\AttributeTok{buzzard=}\DecValTok{52}\NormalTok{)}
\NormalTok{wing\_length}
\end{Highlighting}
\end{Shaded}

\begin{verbatim}
##     goshawk sparrowhawk     buzzard 
##          58          32          52
\end{verbatim}

\begin{Shaded}
\begin{Highlighting}[]
\NormalTok{wing\_length[}\StringTok{"goshawk"}\NormalTok{]              }\CommentTok{\# Access by name: 58}
\end{Highlighting}
\end{Shaded}

\begin{verbatim}
## goshawk 
##      58
\end{verbatim}

\begin{Shaded}
\begin{Highlighting}[]
\NormalTok{wing\_length[}\FunctionTok{c}\NormalTok{(}\StringTok{"goshawk"}\NormalTok{, }\StringTok{"buzzard"}\NormalTok{)] }\CommentTok{\# Multiple names}
\end{Highlighting}
\end{Shaded}

\begin{verbatim}
## goshawk buzzard 
##      58      52
\end{verbatim}

\subsection{Logical Indexing
(Filtering)}\label{logical-indexing-filtering}

Select elements based on conditions:

\begin{Shaded}
\begin{Highlighting}[]
\NormalTok{body\_mass                   }\CommentTok{\# 1200, 1350, 1150, 1425, 1275}
\end{Highlighting}
\end{Shaded}

\begin{verbatim}
## [1] 1200 1350 1150 1425 1275
\end{verbatim}

\begin{Shaded}
\begin{Highlighting}[]
\NormalTok{body\_mass }\SpecialCharTok{\textgreater{}} \DecValTok{1300}            \CommentTok{\# FALSE, TRUE, FALSE, TRUE, FALSE}
\end{Highlighting}
\end{Shaded}

\begin{verbatim}
## [1] FALSE  TRUE FALSE  TRUE FALSE
\end{verbatim}

\begin{Shaded}
\begin{Highlighting}[]
\NormalTok{body\_mass[body\_mass }\SpecialCharTok{\textgreater{}} \DecValTok{1300}\NormalTok{] }\CommentTok{\# Only values \textgreater{} 1300: 1350, 1425}
\end{Highlighting}
\end{Shaded}

\begin{verbatim}
## [1] 1350 1425
\end{verbatim}

\begin{Shaded}
\begin{Highlighting}[]
\CommentTok{\# Which birds are above average weight?}
\NormalTok{mean\_mass }\OtherTok{\textless{}{-}} \FunctionTok{mean}\NormalTok{(body\_mass)}
\NormalTok{body\_mass[body\_mass }\SpecialCharTok{\textgreater{}}\NormalTok{ mean\_mass]}
\end{Highlighting}
\end{Shaded}

\begin{verbatim}
## [1] 1350 1425
\end{verbatim}

\subsection{Vectorized Operations}\label{vectorized-operations}

Operations apply to ALL elements automatically:

\begin{Shaded}
\begin{Highlighting}[]
\NormalTok{body\_mass\_kg }\OtherTok{\textless{}{-}}\NormalTok{ body\_mass }\SpecialCharTok{/} \DecValTok{1000}     \CommentTok{\# Convert all to kg}
\NormalTok{body\_mass\_kg}
\end{Highlighting}
\end{Shaded}

\begin{verbatim}
## [1] 1.200 1.350 1.150 1.425 1.275
\end{verbatim}

\begin{Shaded}
\begin{Highlighting}[]
\CommentTok{\# Element{-}wise operations between vectors:}
\NormalTok{length\_cm }\OtherTok{\textless{}{-}} \FunctionTok{c}\NormalTok{(}\DecValTok{55}\NormalTok{, }\DecValTok{35}\NormalTok{, }\DecValTok{50}\NormalTok{, }\DecValTok{30}\NormalTok{, }\DecValTok{45}\NormalTok{)}
\NormalTok{width\_cm }\OtherTok{\textless{}{-}} \FunctionTok{c}\NormalTok{(}\DecValTok{8}\NormalTok{, }\DecValTok{5}\NormalTok{, }\DecValTok{7}\NormalTok{, }\DecValTok{4}\NormalTok{, }\DecValTok{6}\NormalTok{)}
\NormalTok{area }\OtherTok{\textless{}{-}}\NormalTok{ length\_cm }\SpecialCharTok{*}\NormalTok{ width\_cm         }\CommentTok{\# Multiply corresponding elements}
\end{Highlighting}
\end{Shaded}

\subsubsection{EXERCISE}\label{exercise-3}

\begin{enumerate}
\def\labelenumi{\arabic{enumi}.}
\tightlist
\item
  Create a vector called \texttt{temps} with these daily temperatures:
  18.5, 20.2, 19.8, 22.1, 21.5, 19.0, 17.8
\item
  Calculate the mean temperature
\item
  Find which days were above 20 degrees
\item
  What was the temperature range (max - min)?
\end{enumerate}

\begin{Shaded}
\begin{Highlighting}[]
\CommentTok{\# Your code here:}
\NormalTok{temps }\OtherTok{\textless{}{-}} \FunctionTok{c}\NormalTok{( }\FloatTok{18.5}\NormalTok{, }\FloatTok{20.2}\NormalTok{, }\FloatTok{19.8}\NormalTok{, }\FloatTok{22.1}\NormalTok{, }\FloatTok{21.5}\NormalTok{, }\FloatTok{19.0}\NormalTok{, }\FloatTok{17.8}\NormalTok{)}
\FunctionTok{mean}\NormalTok{(temps)}
\end{Highlighting}
\end{Shaded}

\begin{verbatim}
## [1] 19.84286
\end{verbatim}

\begin{Shaded}
\begin{Highlighting}[]
\FunctionTok{which}\NormalTok{(temps }\SpecialCharTok{\textgreater{}} \DecValTok{20}\NormalTok{)}
\end{Highlighting}
\end{Shaded}

\begin{verbatim}
## [1] 2 4 5
\end{verbatim}

\begin{center}\rule{0.5\linewidth}{0.5pt}\end{center}

\section{PART 6: Data Frames}\label{part-6-data-frames}

A \textbf{data frame} is like a spreadsheet:

\begin{itemize}
\tightlist
\item
  Rows are observations (e.g., individual birds)
\item
  Columns are variables (e.g., species, weight, wing\_length)
\item
  Each column is a vector (so all values in a column are the same type)
\item
  Different columns can have different types
\end{itemize}

\subsection{Creating a Data Frame}\label{creating-a-data-frame}

\begin{Shaded}
\begin{Highlighting}[]
\NormalTok{birds }\OtherTok{\textless{}{-}} \FunctionTok{data.frame}\NormalTok{(}
  \AttributeTok{species =} \FunctionTok{c}\NormalTok{(}\StringTok{"Goshawk"}\NormalTok{, }\StringTok{"Sparrowhawk"}\NormalTok{, }\StringTok{"Buzzard"}\NormalTok{, }\StringTok{"Kestrel"}\NormalTok{, }\StringTok{"Peregrine"}\NormalTok{),}
  \AttributeTok{body\_mass\_g =} \FunctionTok{c}\NormalTok{(}\DecValTok{1200}\NormalTok{, }\DecValTok{250}\NormalTok{, }\DecValTok{900}\NormalTok{, }\DecValTok{200}\NormalTok{, }\DecValTok{750}\NormalTok{),}
  \AttributeTok{wing\_span\_cm =} \FunctionTok{c}\NormalTok{(}\DecValTok{115}\NormalTok{, }\DecValTok{65}\NormalTok{, }\DecValTok{120}\NormalTok{, }\DecValTok{75}\NormalTok{, }\DecValTok{100}\NormalTok{),}
  \AttributeTok{is\_migratory =} \FunctionTok{c}\NormalTok{(}\ConstantTok{FALSE}\NormalTok{, }\ConstantTok{FALSE}\NormalTok{, }\ConstantTok{TRUE}\NormalTok{, }\ConstantTok{TRUE}\NormalTok{, }\ConstantTok{TRUE}\NormalTok{)}
\NormalTok{)}

\CommentTok{\# View the data frame:}
\NormalTok{birds}
\end{Highlighting}
\end{Shaded}

\begin{verbatim}
##       species body_mass_g wing_span_cm is_migratory
## 1     Goshawk        1200          115        FALSE
## 2 Sparrowhawk         250           65        FALSE
## 3     Buzzard         900          120         TRUE
## 4     Kestrel         200           75         TRUE
## 5   Peregrine         750          100         TRUE
\end{verbatim}

For large datasets, use:

\begin{Shaded}
\begin{Highlighting}[]
\FunctionTok{head}\NormalTok{(birds, }\DecValTok{3}\NormalTok{)      }\CommentTok{\# First 3 rows}
\end{Highlighting}
\end{Shaded}

\begin{verbatim}
##       species body_mass_g wing_span_cm is_migratory
## 1     Goshawk        1200          115        FALSE
## 2 Sparrowhawk         250           65        FALSE
## 3     Buzzard         900          120         TRUE
\end{verbatim}

\begin{Shaded}
\begin{Highlighting}[]
\FunctionTok{tail}\NormalTok{(birds, }\DecValTok{2}\NormalTok{)      }\CommentTok{\# Last 2 rows}
\end{Highlighting}
\end{Shaded}

\begin{verbatim}
##     species body_mass_g wing_span_cm is_migratory
## 4   Kestrel         200           75         TRUE
## 5 Peregrine         750          100         TRUE
\end{verbatim}

\subsection{Exploring Data Frame
Structure}\label{exploring-data-frame-structure}

\begin{Shaded}
\begin{Highlighting}[]
\FunctionTok{nrow}\NormalTok{(birds)         }\CommentTok{\# Number of rows: 5}
\end{Highlighting}
\end{Shaded}

\begin{verbatim}
## [1] 5
\end{verbatim}

\begin{Shaded}
\begin{Highlighting}[]
\FunctionTok{ncol}\NormalTok{(birds)         }\CommentTok{\# Number of columns: 4}
\end{Highlighting}
\end{Shaded}

\begin{verbatim}
## [1] 4
\end{verbatim}

\begin{Shaded}
\begin{Highlighting}[]
\FunctionTok{dim}\NormalTok{(birds)          }\CommentTok{\# Both: 5 4}
\end{Highlighting}
\end{Shaded}

\begin{verbatim}
## [1] 5 4
\end{verbatim}

\begin{Shaded}
\begin{Highlighting}[]
\FunctionTok{names}\NormalTok{(birds)        }\CommentTok{\# Column names}
\end{Highlighting}
\end{Shaded}

\begin{verbatim}
## [1] "species"      "body_mass_g"  "wing_span_cm" "is_migratory"
\end{verbatim}

\begin{Shaded}
\begin{Highlighting}[]
\FunctionTok{colnames}\NormalTok{(birds)     }\CommentTok{\# Same as names()}
\end{Highlighting}
\end{Shaded}

\begin{verbatim}
## [1] "species"      "body_mass_g"  "wing_span_cm" "is_migratory"
\end{verbatim}

\begin{Shaded}
\begin{Highlighting}[]
\FunctionTok{str}\NormalTok{(birds)          }\CommentTok{\# Structure {-} shows types of each column}
\end{Highlighting}
\end{Shaded}

\begin{verbatim}
## 'data.frame':    5 obs. of  4 variables:
##  $ species     : chr  "Goshawk" "Sparrowhawk" "Buzzard" "Kestrel" ...
##  $ body_mass_g : num  1200 250 900 200 750
##  $ wing_span_cm: num  115 65 120 75 100
##  $ is_migratory: logi  FALSE FALSE TRUE TRUE TRUE
\end{verbatim}

\begin{Shaded}
\begin{Highlighting}[]
\FunctionTok{summary}\NormalTok{(birds)      }\CommentTok{\# Statistical summary of each column}
\end{Highlighting}
\end{Shaded}

\begin{verbatim}
##    species           body_mass_g    wing_span_cm is_migratory   
##  Length:5           Min.   : 200   Min.   : 65   Mode :logical  
##  Class :character   1st Qu.: 250   1st Qu.: 75   FALSE:2        
##  Mode  :character   Median : 750   Median :100   TRUE :3        
##                     Mean   : 660   Mean   : 95                  
##                     3rd Qu.: 900   3rd Qu.:115                  
##                     Max.   :1200   Max.   :120
\end{verbatim}

\subsection{Accessing Columns}\label{accessing-columns}

\begin{Shaded}
\begin{Highlighting}[]
\CommentTok{\# Method 1: $ notation (most common)}
\NormalTok{birds}\SpecialCharTok{$}\NormalTok{species           }\CommentTok{\# Returns species column as a vector}
\end{Highlighting}
\end{Shaded}

\begin{verbatim}
## [1] "Goshawk"     "Sparrowhawk" "Buzzard"     "Kestrel"     "Peregrine"
\end{verbatim}

\begin{Shaded}
\begin{Highlighting}[]
\CommentTok{\#$ references a column in a dataframe}
\NormalTok{birds}\SpecialCharTok{$}\NormalTok{body\_mass\_g      }\CommentTok{\# Returns body\_mass\_g column}
\end{Highlighting}
\end{Shaded}

\begin{verbatim}
## [1] 1200  250  900  200  750
\end{verbatim}

\begin{Shaded}
\begin{Highlighting}[]
\CommentTok{\# Method 2: Double brackets}
\NormalTok{birds[[}\StringTok{"species"}\NormalTok{]]      }\CommentTok{\# Same result}
\end{Highlighting}
\end{Shaded}

\begin{verbatim}
## [1] "Goshawk"     "Sparrowhawk" "Buzzard"     "Kestrel"     "Peregrine"
\end{verbatim}

\begin{Shaded}
\begin{Highlighting}[]
\CommentTok{\# Method 3: Single brackets (returns data frame)}
\NormalTok{birds[, }\StringTok{"species"}\NormalTok{]      }\CommentTok{\# All rows, species column}
\end{Highlighting}
\end{Shaded}

\begin{verbatim}
## [1] "Goshawk"     "Sparrowhawk" "Buzzard"     "Kestrel"     "Peregrine"
\end{verbatim}

\begin{Shaded}
\begin{Highlighting}[]
\NormalTok{birds[, }\FunctionTok{c}\NormalTok{(}\StringTok{"species"}\NormalTok{, }\StringTok{"body\_mass\_g"}\NormalTok{)]  }\CommentTok{\# Multiple columns}
\end{Highlighting}
\end{Shaded}

\begin{verbatim}
##       species body_mass_g
## 1     Goshawk        1200
## 2 Sparrowhawk         250
## 3     Buzzard         900
## 4     Kestrel         200
## 5   Peregrine         750
\end{verbatim}

\subsection{Accessing Rows}\label{accessing-rows}

\begin{Shaded}
\begin{Highlighting}[]
\NormalTok{birds[}\DecValTok{1}\NormalTok{, ]              }\CommentTok{\# First row (all columns)}
\end{Highlighting}
\end{Shaded}

\begin{verbatim}
##   species body_mass_g wing_span_cm is_migratory
## 1 Goshawk        1200          115        FALSE
\end{verbatim}

\begin{Shaded}
\begin{Highlighting}[]
\NormalTok{birds[}\FunctionTok{c}\NormalTok{(}\DecValTok{1}\NormalTok{, }\DecValTok{3}\NormalTok{), ]        }\CommentTok{\# Rows 1 and 3}
\end{Highlighting}
\end{Shaded}

\begin{verbatim}
##   species body_mass_g wing_span_cm is_migratory
## 1 Goshawk        1200          115        FALSE
## 3 Buzzard         900          120         TRUE
\end{verbatim}

\begin{Shaded}
\begin{Highlighting}[]
\NormalTok{birds[}\DecValTok{1}\SpecialCharTok{:}\DecValTok{3}\NormalTok{, ]            }\CommentTok{\# Rows 1 through 3}
\end{Highlighting}
\end{Shaded}

\begin{verbatim}
##       species body_mass_g wing_span_cm is_migratory
## 1     Goshawk        1200          115        FALSE
## 2 Sparrowhawk         250           65        FALSE
## 3     Buzzard         900          120         TRUE
\end{verbatim}

\subsection{Accessing Specific Cells}\label{accessing-specific-cells}

\begin{Shaded}
\begin{Highlighting}[]
\NormalTok{birds[}\DecValTok{1}\NormalTok{, }\StringTok{"species"}\NormalTok{]              }\CommentTok{\# Row 1, species column: "Goshawk"}
\end{Highlighting}
\end{Shaded}

\begin{verbatim}
## [1] "Goshawk"
\end{verbatim}

\begin{Shaded}
\begin{Highlighting}[]
\NormalTok{birds[}\DecValTok{2}\NormalTok{, }\DecValTok{2}\NormalTok{]                      }\CommentTok{\# Row 2, column 2: 250}
\end{Highlighting}
\end{Shaded}

\begin{verbatim}
## [1] 250
\end{verbatim}

\subsection{Filtering Rows with
Conditions}\label{filtering-rows-with-conditions}

\begin{Shaded}
\begin{Highlighting}[]
\CommentTok{\# Birds heavier than 500g:}
\NormalTok{birds[birds}\SpecialCharTok{$}\NormalTok{body\_mass\_g }\SpecialCharTok{\textgreater{}} \DecValTok{500}\NormalTok{, ]}
\end{Highlighting}
\end{Shaded}

\begin{verbatim}
##     species body_mass_g wing_span_cm is_migratory
## 1   Goshawk        1200          115        FALSE
## 3   Buzzard         900          120         TRUE
## 5 Peregrine         750          100         TRUE
\end{verbatim}

\begin{Shaded}
\begin{Highlighting}[]
\CommentTok{\# Migratory birds only:}
\NormalTok{birds[birds}\SpecialCharTok{$}\NormalTok{is\_migratory }\SpecialCharTok{==} \ConstantTok{TRUE}\NormalTok{, ]}
\end{Highlighting}
\end{Shaded}

\begin{verbatim}
##     species body_mass_g wing_span_cm is_migratory
## 3   Buzzard         900          120         TRUE
## 4   Kestrel         200           75         TRUE
## 5 Peregrine         750          100         TRUE
\end{verbatim}

\begin{Shaded}
\begin{Highlighting}[]
\CommentTok{\# Or equivalently:}
\NormalTok{birds[birds}\SpecialCharTok{$}\NormalTok{is\_migratory, ]}
\end{Highlighting}
\end{Shaded}

\begin{verbatim}
##     species body_mass_g wing_span_cm is_migratory
## 3   Buzzard         900          120         TRUE
## 4   Kestrel         200           75         TRUE
## 5 Peregrine         750          100         TRUE
\end{verbatim}

\begin{Shaded}
\begin{Highlighting}[]
\CommentTok{\# Multiple conditions (use \& for AND, | for OR):}
\NormalTok{birds[birds}\SpecialCharTok{$}\NormalTok{body\_mass\_g }\SpecialCharTok{\textgreater{}} \DecValTok{500} \SpecialCharTok{\&}\NormalTok{ birds}\SpecialCharTok{$}\NormalTok{is\_migratory }\SpecialCharTok{==} \ConstantTok{TRUE}\NormalTok{, ]}
\end{Highlighting}
\end{Shaded}

\begin{verbatim}
##     species body_mass_g wing_span_cm is_migratory
## 3   Buzzard         900          120         TRUE
## 5 Peregrine         750          100         TRUE
\end{verbatim}

\subsection{Adding New Columns}\label{adding-new-columns}

\begin{Shaded}
\begin{Highlighting}[]
\CommentTok{\# Calculate body mass in kg:}
\NormalTok{birds}\SpecialCharTok{$}\NormalTok{body\_mass\_kg }\OtherTok{\textless{}{-}}\NormalTok{ birds}\SpecialCharTok{$}\NormalTok{body\_mass\_g }\SpecialCharTok{/} \DecValTok{1000}
\NormalTok{birds}
\end{Highlighting}
\end{Shaded}

\begin{verbatim}
##       species body_mass_g wing_span_cm is_migratory body_mass_kg
## 1     Goshawk        1200          115        FALSE         1.20
## 2 Sparrowhawk         250           65        FALSE         0.25
## 3     Buzzard         900          120         TRUE         0.90
## 4     Kestrel         200           75         TRUE         0.20
## 5   Peregrine         750          100         TRUE         0.75
\end{verbatim}

\begin{Shaded}
\begin{Highlighting}[]
\CommentTok{\# Add a categorical column:}
\NormalTok{birds}\SpecialCharTok{$}\NormalTok{size\_category }\OtherTok{\textless{}{-}} \FunctionTok{ifelse}\NormalTok{(birds}\SpecialCharTok{$}\NormalTok{body\_mass\_g }\SpecialCharTok{\textgreater{}} \DecValTok{500}\NormalTok{, }\StringTok{"Large"}\NormalTok{, }\StringTok{"Small"}\NormalTok{) }\CommentTok{\#ifelse(test, yes, no)}
\NormalTok{birds}
\end{Highlighting}
\end{Shaded}

\begin{verbatim}
##       species body_mass_g wing_span_cm is_migratory body_mass_kg size_category
## 1     Goshawk        1200          115        FALSE         1.20         Large
## 2 Sparrowhawk         250           65        FALSE         0.25         Small
## 3     Buzzard         900          120         TRUE         0.90         Large
## 4     Kestrel         200           75         TRUE         0.20         Small
## 5   Peregrine         750          100         TRUE         0.75         Large
\end{verbatim}

\subsection{Modifying Data}\label{modifying-data}

\begin{Shaded}
\begin{Highlighting}[]
\CommentTok{\# Change a single value:}
\NormalTok{birds[}\DecValTok{1}\NormalTok{, }\StringTok{"body\_mass\_g"}\NormalTok{] }\OtherTok{\textless{}{-}} \DecValTok{1250}
\NormalTok{birds}
\end{Highlighting}
\end{Shaded}

\begin{verbatim}
##       species body_mass_g wing_span_cm is_migratory body_mass_kg size_category
## 1     Goshawk        1250          115        FALSE         1.20         Large
## 2 Sparrowhawk         250           65        FALSE         0.25         Small
## 3     Buzzard         900          120         TRUE         0.90         Large
## 4     Kestrel         200           75         TRUE         0.20         Small
## 5   Peregrine         750          100         TRUE         0.75         Large
\end{verbatim}

\begin{Shaded}
\begin{Highlighting}[]
\CommentTok{\# Update based on condition:}
\NormalTok{birds}\SpecialCharTok{$}\NormalTok{body\_mass\_g[birds}\SpecialCharTok{$}\NormalTok{species }\SpecialCharTok{==} \StringTok{"Goshawk"}\NormalTok{] }\OtherTok{\textless{}{-}} \DecValTok{1300}
\end{Highlighting}
\end{Shaded}

\subsection{Sorting Data Frames}\label{sorting-data-frames}

\begin{Shaded}
\begin{Highlighting}[]
\CommentTok{\# Sort by body mass (ascending):}
\NormalTok{birds[}\FunctionTok{order}\NormalTok{(birds}\SpecialCharTok{$}\NormalTok{body\_mass\_g), ]}
\end{Highlighting}
\end{Shaded}

\begin{verbatim}
##       species body_mass_g wing_span_cm is_migratory body_mass_kg size_category
## 4     Kestrel         200           75         TRUE         0.20         Small
## 2 Sparrowhawk         250           65        FALSE         0.25         Small
## 5   Peregrine         750          100         TRUE         0.75         Large
## 3     Buzzard         900          120         TRUE         0.90         Large
## 1     Goshawk        1300          115        FALSE         1.20         Large
\end{verbatim}

\begin{Shaded}
\begin{Highlighting}[]
\CommentTok{\# Sort by body mass (descending):}
\NormalTok{birds[}\FunctionTok{order}\NormalTok{(birds}\SpecialCharTok{$}\NormalTok{body\_mass\_g, }\AttributeTok{decreasing =} \ConstantTok{TRUE}\NormalTok{), ]}
\end{Highlighting}
\end{Shaded}

\begin{verbatim}
##       species body_mass_g wing_span_cm is_migratory body_mass_kg size_category
## 1     Goshawk        1300          115        FALSE         1.20         Large
## 3     Buzzard         900          120         TRUE         0.90         Large
## 5   Peregrine         750          100         TRUE         0.75         Large
## 2 Sparrowhawk         250           65        FALSE         0.25         Small
## 4     Kestrel         200           75         TRUE         0.20         Small
\end{verbatim}

\begin{Shaded}
\begin{Highlighting}[]
\CommentTok{\# Or use negative for numeric columns:}
\NormalTok{birds[}\FunctionTok{order}\NormalTok{(}\SpecialCharTok{{-}}\NormalTok{birds}\SpecialCharTok{$}\NormalTok{body\_mass\_g), ]}
\end{Highlighting}
\end{Shaded}

\begin{verbatim}
##       species body_mass_g wing_span_cm is_migratory body_mass_kg size_category
## 1     Goshawk        1300          115        FALSE         1.20         Large
## 3     Buzzard         900          120         TRUE         0.90         Large
## 5   Peregrine         750          100         TRUE         0.75         Large
## 2 Sparrowhawk         250           65        FALSE         0.25         Small
## 4     Kestrel         200           75         TRUE         0.20         Small
\end{verbatim}

\subsubsection{EXERCISE}\label{exercise-4}

\begin{enumerate}
\def\labelenumi{\arabic{enumi}.}
\tightlist
\item
  Create a data frame called \texttt{trees} with columns:

  \begin{itemize}
  \tightlist
  \item
    species: ``Oak'', ``Maple'', ``Pine'', ``Birch''
  \item
    height\_m: 25, 20, 30, 18
  \item
    is\_evergreen: FALSE, FALSE, TRUE, FALSE
  \end{itemize}
\item
  Add a new column \texttt{height\_ft} (height in feet = height\_m *
  3.28)
\item
  Filter to show only trees taller than 20 meters
\item
  What is the average height of all trees in meters?
\end{enumerate}

\begin{Shaded}
\begin{Highlighting}[]
\CommentTok{\# Your code here:}
\NormalTok{trees }\OtherTok{\textless{}{-}} \FunctionTok{data.frame}\NormalTok{(}\AttributeTok{species =} \FunctionTok{c}\NormalTok{(}\StringTok{"oak"}\NormalTok{, }\StringTok{"maple"}\NormalTok{, }\StringTok{"pine"}\NormalTok{, }\StringTok{"birch"}\NormalTok{), }\AttributeTok{height\_m =}\FunctionTok{c}\NormalTok{(}\DecValTok{25}\NormalTok{, }\DecValTok{20}\NormalTok{, }\DecValTok{30}\NormalTok{, }\DecValTok{18}\NormalTok{), }\AttributeTok{is\_evergreen =} \FunctionTok{c}\NormalTok{(}\ConstantTok{FALSE}\NormalTok{, }\ConstantTok{FALSE}\NormalTok{, }\ConstantTok{TRUE}\NormalTok{, }\ConstantTok{FALSE}\NormalTok{)}
\NormalTok{)}

\NormalTok{trees}\SpecialCharTok{$}\NormalTok{height\_ft }\OtherTok{\textless{}{-}}\NormalTok{ trees}\SpecialCharTok{$}\NormalTok{height\_m}\SpecialCharTok{*}\FloatTok{3.28}
\FunctionTok{which}\NormalTok{(trees}\SpecialCharTok{$}\NormalTok{height\_m }\SpecialCharTok{\textgreater{}}\DecValTok{20}\NormalTok{)}
\end{Highlighting}
\end{Shaded}

\begin{verbatim}
## [1] 1 3
\end{verbatim}

\begin{center}\rule{0.5\linewidth}{0.5pt}\end{center}

\section{PART 7: Handling Missing Data (NA
Values)}\label{part-7-handling-missing-data-na-values}

Real-world biological data almost ALWAYS has missing values, in R,
missing values are represented as NA.

\begin{Shaded}
\begin{Highlighting}[]
\CommentTok{\# Create data with missing values:}
\NormalTok{measurements }\OtherTok{\textless{}{-}} \FunctionTok{data.frame}\NormalTok{(}
  \AttributeTok{sample\_id =} \DecValTok{1}\SpecialCharTok{:}\DecValTok{6}\NormalTok{,}
  \AttributeTok{weight =} \FunctionTok{c}\NormalTok{(}\FloatTok{45.2}\NormalTok{, }\FloatTok{52.1}\NormalTok{, }\ConstantTok{NA}\NormalTok{, }\FloatTok{48.7}\NormalTok{, }\FloatTok{51.3}\NormalTok{, }\ConstantTok{NA}\NormalTok{),}
  \AttributeTok{length =} \FunctionTok{c}\NormalTok{(}\FloatTok{12.5}\NormalTok{, }\ConstantTok{NA}\NormalTok{, }\FloatTok{14.2}\NormalTok{, }\FloatTok{13.8}\NormalTok{, }\ConstantTok{NA}\NormalTok{, }\FloatTok{15.1}\NormalTok{),}
  \AttributeTok{status =} \FunctionTok{c}\NormalTok{(}\StringTok{"alive"}\NormalTok{, }\StringTok{"alive"}\NormalTok{, }\StringTok{"dead"}\NormalTok{, }\ConstantTok{NA}\NormalTok{, }\StringTok{"alive"}\NormalTok{, }\StringTok{"dead"}\NormalTok{)}
\NormalTok{)}
\NormalTok{measurements}
\end{Highlighting}
\end{Shaded}

\begin{verbatim}
##   sample_id weight length status
## 1         1   45.2   12.5  alive
## 2         2   52.1     NA  alive
## 3         3     NA   14.2   dead
## 4         4   48.7   13.8   <NA>
## 5         5   51.3     NA  alive
## 6         6     NA   15.1   dead
\end{verbatim}

\subsection{Detecting Missing Values}\label{detecting-missing-values}

\begin{Shaded}
\begin{Highlighting}[]
\FunctionTok{is.na}\NormalTok{(measurements}\SpecialCharTok{$}\NormalTok{weight)           }\CommentTok{\# TRUE/FALSE for each element}
\end{Highlighting}
\end{Shaded}

\begin{verbatim}
## [1] FALSE FALSE  TRUE FALSE FALSE  TRUE
\end{verbatim}

\begin{Shaded}
\begin{Highlighting}[]
\FunctionTok{sum}\NormalTok{(}\FunctionTok{is.na}\NormalTok{(measurements}\SpecialCharTok{$}\NormalTok{weight))      }\CommentTok{\# Count NAs in weight: 2}
\end{Highlighting}
\end{Shaded}

\begin{verbatim}
## [1] 2
\end{verbatim}

\begin{Shaded}
\begin{Highlighting}[]
\FunctionTok{sum}\NormalTok{(}\FunctionTok{is.na}\NormalTok{(measurements))             }\CommentTok{\# Total NAs in whole data frame}
\end{Highlighting}
\end{Shaded}

\begin{verbatim}
## [1] 5
\end{verbatim}

\begin{Shaded}
\begin{Highlighting}[]
\CommentTok{\# Which rows have complete data?}
\FunctionTok{complete.cases}\NormalTok{(measurements)          }\CommentTok{\# TRUE/FALSE for each row}
\end{Highlighting}
\end{Shaded}

\begin{verbatim}
## [1]  TRUE FALSE FALSE FALSE FALSE FALSE
\end{verbatim}

\subsection{NA in Calculations}\label{na-in-calculations}

\begin{Shaded}
\begin{Highlighting}[]
\FunctionTok{mean}\NormalTok{(measurements}\SpecialCharTok{$}\NormalTok{weight)             }
\end{Highlighting}
\end{Shaded}

\begin{verbatim}
## [1] NA
\end{verbatim}

\begin{Shaded}
\begin{Highlighting}[]
\CommentTok{\# Solution: use na.rm = TRUE (remove NAs)}
\FunctionTok{mean}\NormalTok{(measurements}\SpecialCharTok{$}\NormalTok{weight, }\AttributeTok{na.rm =} \ConstantTok{TRUE}\NormalTok{)   }
\end{Highlighting}
\end{Shaded}

\begin{verbatim}
## [1] 49.325
\end{verbatim}

\begin{Shaded}
\begin{Highlighting}[]
\CommentTok{\# This works for most statistical functions:}
\FunctionTok{sum}\NormalTok{(measurements}\SpecialCharTok{$}\NormalTok{weight, }\AttributeTok{na.rm =} \ConstantTok{TRUE}\NormalTok{)}
\end{Highlighting}
\end{Shaded}

\begin{verbatim}
## [1] 197.3
\end{verbatim}

\begin{Shaded}
\begin{Highlighting}[]
\FunctionTok{sd}\NormalTok{(measurements}\SpecialCharTok{$}\NormalTok{length, }\AttributeTok{na.rm =} \ConstantTok{TRUE}\NormalTok{)}
\end{Highlighting}
\end{Shaded}

\begin{verbatim}
## [1] 1.080123
\end{verbatim}

\begin{Shaded}
\begin{Highlighting}[]
\FunctionTok{max}\NormalTok{(measurements}\SpecialCharTok{$}\NormalTok{weight, }\AttributeTok{na.rm =} \ConstantTok{TRUE}\NormalTok{)}
\end{Highlighting}
\end{Shaded}

\begin{verbatim}
## [1] 52.1
\end{verbatim}

\subsection{Removing Rows with Missing
Values}\label{removing-rows-with-missing-values}

\begin{Shaded}
\begin{Highlighting}[]
\CommentTok{\# Remove rows with ANY missing values:}
\NormalTok{measurements\_complete }\OtherTok{\textless{}{-}} \FunctionTok{na.omit}\NormalTok{(measurements)}
\NormalTok{measurements\_complete   }
\end{Highlighting}
\end{Shaded}

\begin{verbatim}
##   sample_id weight length status
## 1         1   45.2   12.5  alive
\end{verbatim}

\begin{Shaded}
\begin{Highlighting}[]
\CommentTok{\# Remove rows with NA in a specific column:}
\NormalTok{measurements[}\SpecialCharTok{!}\FunctionTok{is.na}\NormalTok{(measurements}\SpecialCharTok{$}\NormalTok{weight), ]}
\end{Highlighting}
\end{Shaded}

\begin{verbatim}
##   sample_id weight length status
## 1         1   45.2   12.5  alive
## 2         2   52.1     NA  alive
## 4         4   48.7   13.8   <NA>
## 5         5   51.3     NA  alive
\end{verbatim}

\subsection{Replacing Missing Values}\label{replacing-missing-values}

\begin{Shaded}
\begin{Highlighting}[]
\CommentTok{\# Replace NAs with a specific value:}
\NormalTok{measurements}\SpecialCharTok{$}\NormalTok{weight[}\FunctionTok{is.na}\NormalTok{(measurements}\SpecialCharTok{$}\NormalTok{weight)] }\OtherTok{\textless{}{-}} \DecValTok{0}

\CommentTok{\# Replace with mean:}
\CommentTok{\# First, let\textquotesingle{}s reset the data}
\NormalTok{measurements}\SpecialCharTok{$}\NormalTok{weight }\OtherTok{\textless{}{-}} \FunctionTok{c}\NormalTok{(}\FloatTok{45.2}\NormalTok{, }\FloatTok{52.1}\NormalTok{, }\ConstantTok{NA}\NormalTok{, }\FloatTok{48.7}\NormalTok{, }\FloatTok{51.3}\NormalTok{, }\ConstantTok{NA}\NormalTok{)}

\CommentTok{\# Replace NAs with the mean weight}
\NormalTok{mean\_weight }\OtherTok{\textless{}{-}} \FunctionTok{mean}\NormalTok{(measurements}\SpecialCharTok{$}\NormalTok{weight, }\AttributeTok{na.rm =} \ConstantTok{TRUE}\NormalTok{)}
\NormalTok{measurements}\SpecialCharTok{$}\NormalTok{weight[}\FunctionTok{is.na}\NormalTok{(measurements}\SpecialCharTok{$}\NormalTok{weight)] }\OtherTok{\textless{}{-}}\NormalTok{ mean\_weight}
\NormalTok{measurements}
\end{Highlighting}
\end{Shaded}

\begin{verbatim}
##   sample_id weight length status
## 1         1 45.200   12.5  alive
## 2         2 52.100     NA  alive
## 3         3 49.325   14.2   dead
## 4         4 48.700   13.8   <NA>
## 5         5 51.300     NA  alive
## 6         6 49.325   15.1   dead
\end{verbatim}

\subsubsection{EXERCISE}\label{exercise-5}

\begin{enumerate}
\def\labelenumi{\arabic{enumi}.}
\tightlist
\item
  How many missing values are in the length column of measurements?
\item
  What is the median length (excluding NAs)?
\item
  Remove all rows with any missing values from measurements
\end{enumerate}

\begin{Shaded}
\begin{Highlighting}[]
\CommentTok{\# Your code here:}
\end{Highlighting}
\end{Shaded}

\begin{center}\rule{0.5\linewidth}{0.5pt}\end{center}

\section{PART 8: Packages - Extending R's
Capabilities}\label{part-8-packages---extending-rs-capabilities}

R comes with many built-in functions, \textbf{packages} add many many
more. Packages are collections of functions that other people have
written made avialable to us.

\subsection{Installing Packages}\label{installing-packages}

You only need to install a package ONCE on your computer. Do this in the
Console, not in your script:

\begin{Shaded}
\begin{Highlighting}[]
\FunctionTok{install.packages}\NormalTok{(}\StringTok{"ggplot2"}\NormalTok{)  }\CommentTok{\# For advanced plotting}
\FunctionTok{install.packages}\NormalTok{(}\StringTok{"dplyr"}\NormalTok{)    }\CommentTok{\# For data manipulation}
\end{Highlighting}
\end{Shaded}

\subsection{Loading Packages}\label{loading-packages}

You need to load a package EACH TIME you start R. Put \texttt{library()}
calls at the top of your script:

\begin{Shaded}
\begin{Highlighting}[]
\FunctionTok{library}\NormalTok{(ggplot2)}
\FunctionTok{library}\NormalTok{(dplyr)}
\end{Highlighting}
\end{Shaded}

\subsection{Checking Installed
Packages}\label{checking-installed-packages}

\begin{Shaded}
\begin{Highlighting}[]
\FunctionTok{installed.packages}\NormalTok{()   }\CommentTok{\# Shows all installed packages}
\end{Highlighting}
\end{Shaded}

\section{PART 9: Importing and Exporting
Data}\label{part-9-importing-and-exporting-data}

In real research, you'll read data from files, not type it in!

\subsection{Setting Your Working
Directory}\label{setting-your-working-directory}

\begin{Shaded}
\begin{Highlighting}[]
\CommentTok{\# First, check where R is looking for files:}
\FunctionTok{getwd}\NormalTok{()}
\end{Highlighting}
\end{Shaded}

\begin{verbatim}
## [1] "/Users/tabi/Downloads/4th Year/Biologists Toolkit/BIOL3872/Lab 2/Scripts"
\end{verbatim}

\begin{Shaded}
\begin{Highlighting}[]
\CommentTok{\# Change the working directory:}
\FunctionTok{setwd}\NormalTok{(}\StringTok{"/path/to/your/folder"}\NormalTok{)}
\CommentTok{\# Or use: Session → Set Working Directory → Choose Directory in RStudio}
\end{Highlighting}
\end{Shaded}

\subsection{Reading CSV Files}\label{reading-csv-files}

CSV (Comma-Separated Values) is the most common format.

\begin{Shaded}
\begin{Highlighting}[]
\CommentTok{\# Basic syntax:}
\NormalTok{my\_data }\OtherTok{\textless{}{-}} \FunctionTok{read.csv}\NormalTok{(}\StringTok{"filename.csv"}\NormalTok{)}

\CommentTok{\# With options:}
\NormalTok{my\_data }\OtherTok{\textless{}{-}} \FunctionTok{read.csv}\NormalTok{(}\StringTok{"filename.csv"}\NormalTok{,}
                    \AttributeTok{header =} \ConstantTok{TRUE}\NormalTok{,              }\CommentTok{\# First row is column names}
                    \AttributeTok{stringsAsFactors =} \ConstantTok{FALSE}\NormalTok{,   }\CommentTok{\# Keep text as text}
                    \AttributeTok{na.strings =} \FunctionTok{c}\NormalTok{(}\StringTok{""}\NormalTok{, }\StringTok{"NA"}\NormalTok{, }\StringTok{"N/A"}\NormalTok{))  }\CommentTok{\# What counts as NA}
\end{Highlighting}
\end{Shaded}

Let's create a sample file to practice:

\begin{Shaded}
\begin{Highlighting}[]
\NormalTok{sample\_data }\OtherTok{\textless{}{-}} \FunctionTok{data.frame}\NormalTok{(}
  \AttributeTok{species =} \FunctionTok{c}\NormalTok{(}\StringTok{"Goshawk"}\NormalTok{, }\StringTok{"Sparrowhawk"}\NormalTok{, }\StringTok{"Buzzard"}\NormalTok{),}
  \AttributeTok{count =} \FunctionTok{c}\NormalTok{(}\DecValTok{15}\NormalTok{, }\DecValTok{42}\NormalTok{, }\DecValTok{28}\NormalTok{),}
  \AttributeTok{location =} \FunctionTok{c}\NormalTok{(}\StringTok{"Forest"}\NormalTok{, }\StringTok{"Garden"}\NormalTok{, }\StringTok{"Field"}\NormalTok{)}
\NormalTok{)}

\CommentTok{\# Write to CSV:}
\FunctionTok{write.csv}\NormalTok{(sample\_data, }\StringTok{"bird\_counts.csv"}\NormalTok{, }\AttributeTok{row.names =} \ConstantTok{FALSE}\NormalTok{)}

\CommentTok{\# Read it back:}
\NormalTok{bird\_data }\OtherTok{\textless{}{-}} \FunctionTok{read.csv}\NormalTok{(}\StringTok{"bird\_counts.csv"}\NormalTok{)}
\NormalTok{bird\_data}
\end{Highlighting}
\end{Shaded}

\begin{verbatim}
##       species count location
## 1     Goshawk    15   Forest
## 2 Sparrowhawk    42   Garden
## 3     Buzzard    28    Field
\end{verbatim}

\subsection{Always Check Your Data After
Importing}\label{always-check-your-data-after-importing}

\begin{Shaded}
\begin{Highlighting}[]
\FunctionTok{head}\NormalTok{(bird\_data)          }\CommentTok{\# First few rows}
\end{Highlighting}
\end{Shaded}

\begin{verbatim}
##       species count location
## 1     Goshawk    15   Forest
## 2 Sparrowhawk    42   Garden
## 3     Buzzard    28    Field
\end{verbatim}

\begin{Shaded}
\begin{Highlighting}[]
\FunctionTok{str}\NormalTok{(bird\_data)           }\CommentTok{\# Structure}
\end{Highlighting}
\end{Shaded}

\begin{verbatim}
## 'data.frame':    3 obs. of  3 variables:
##  $ species : chr  "Goshawk" "Sparrowhawk" "Buzzard"
##  $ count   : int  15 42 28
##  $ location: chr  "Forest" "Garden" "Field"
\end{verbatim}

\begin{Shaded}
\begin{Highlighting}[]
\FunctionTok{summary}\NormalTok{(bird\_data)       }\CommentTok{\# Summary statistics}
\end{Highlighting}
\end{Shaded}

\begin{verbatim}
##    species              count         location        
##  Length:3           Min.   :15.00   Length:3          
##  Class :character   1st Qu.:21.50   Class :character  
##  Mode  :character   Median :28.00   Mode  :character  
##                     Mean   :28.33                     
##                     3rd Qu.:35.00                     
##                     Max.   :42.00
\end{verbatim}

\begin{Shaded}
\begin{Highlighting}[]
\FunctionTok{dim}\NormalTok{(bird\_data)           }\CommentTok{\# Dimensions}
\end{Highlighting}
\end{Shaded}

\begin{verbatim}
## [1] 3 3
\end{verbatim}

\subsection{Now let's head to the
assignment!}\label{now-lets-head-to-the-assignment}

\end{document}
